Una comunidad tiene inicialmente $100$ habitantes. Su crecimiento interno es proporcional a la población, con tasa numérica $0.15(1+t^2)$, y recibe inmigración neta constante de $20$ personas al año. Use $t$ como valor numérico del tiempo en años y $P$ como población continua del modelo.

La tasa $0.15(1+t^2)$ se expresa en año$^{-1}$ cuando $t$ es el valor numérico en años; de forma dimensional equivale a $0.15\,\mathrm{año}^{-1}[1+(\tau/(1\,\mathrm{año}))^2]$. Se supone crecimiento interno neto con esa tasa, inmigración constante y ausencia de otras entradas o salidas.

\begin{enumerate}[label=(\alph*),ref=\alph*,leftmargin=*]
\item Formule el PVI que describe la población.
\item Calcule el ritmo de crecimiento en el instante inicial.
\item Determine la velocidad de crecimiento cuando la población llega a $500$ habitantes.
\end{enumerate}
